\section{An early--late reduction of the nonlinear regime}
\label{sec:terminal-modal-regime-reduction}

The directed remainder bound can control the global safe envelope once the
position has reached order $q$.  The first few full-face positions are only
of order $q^2$, so the two time ranges must be separated.  This section
proves the late comparison with explicit rational constants and reduces the
early range to one signed half-retention inequality.  The latter inequality
remains open.

Put $p^\star=p_{m+1}^\star$, let $p_0,v_0$ be the literal full-face entry
state, and define
\begin{equation}
 e_0=p_0-p^\star,\qquad f_0=v_0-p^\star,\qquad \delta=1-q.
 \label{eq:terminal-modal-entry-position-errors}
\end{equation}

\begin{lemma}[Elementary full-entry position bounds; proved here]
\label{lem:terminal-modal-entry-position-bounds}
For every $m\ge64$,
\begin{equation}
 \min_jp^\star(j)>\frac{1243}{160}q,
 \qquad
 \max_jp^\star(j)<\frac{157}{20}q,
 \qquad
 e_0\ge-\frac{157}{20}q\one,
 \qquad
 f_0\ge-\frac q8\one.
 \label{eq:terminal-modal-entry-position-bounds}
\end{equation}
\end{lemma}

\begin{proof}
Write
\[
 x=2m\operatorname{artanh}q,
 \qquad t=\chi^m=e^{-x}.
\]
The full optimum in \eqref{eq:terminal-modal-full-optimum} decreases from the
seed to the far endpoint, and its two extreme values are
\begin{equation}
 \frac{p^\star(0)}q=\coth x-\frac15,
 \qquad
 \frac{p^\star(m)}q=\operatorname{csch}x-\frac15.
 \label{eq:terminal-modal-full-optimum-extremes}
\end{equation}
Since $x>1/8$ and
\begin{equation}
 \tanh(1/8)\ge\frac18-\frac{(1/8)^3}{3}
 =\frac{191}{1536}>\frac{20}{161},
 \qquad
 \frac{191}{1536}-\frac{20}{161}=\frac{31}{247296},
 \label{eq:terminal-modal-full-optimum-sharpened-upper}
\end{equation}
the first value is less than $161/20-1/5=157/20$.

For the other direction, $q\le1/1024$ and
\begin{equation}
 x\le\frac1{8(1-q^2)}
 \le x_0:=\frac18+\frac1{4194304}<\frac17.
 \label{eq:terminal-modal-full-optimum-x-upper}
\end{equation}
The elementary series comparison
$\sinh z\le z+z^3/5$ on $0\le z\le1/7$ and the exact rational inequality
\[
 \frac{32}{255}-\left(x_0+\frac{x_0^3}{5}\right)
 =\frac{1868969381144821709}{18815678955183742648320}>0
\]
give $\operatorname{csch}x>255/32$.  This proves the first two bounds.

The proper-prefix chronology gives $p_0\ge0$, so the bound on $e_0$ follows
from the full-optimum maximum.  For $f_0$, let $v_m^{\rm out}$ be the center
after the last proper-prefix step.  On the old rows and the newly admitted
endpoint respectively, exact center transport gives
\begin{equation}
 f_0(j)=v_m^{\rm out}(j)-p_m^\star(j)\quad(j<m),
 \qquad f_0(m)=0.
 \label{eq:terminal-modal-entry-f0-transport}
\end{equation}
The chronology proof also gives $v_m^{\rm out}\ge0$.

It remains to bound the proper-prefix optimum.  Put $T=\chi^m$.  Since
$T\ge7/8>1/5$, one has $b_m<0$, and direct differences in
\eqref{eq:terminal-modal-prefix-optimum} show that $p_m^\star(j)$ decreases
with $j$.  At the seed,
\begin{equation}
 \frac{p_m^\star(0)}q
 =\frac{2(2+T-3T^2)}{5(1+T^2)}.
 \label{eq:terminal-modal-proper-seed-optimum}
\end{equation}
The right side decreases on $[7/8,1]$, because the numerator of its
derivative after clearing positive factors is $1-10T-T^2<0$.
Moreover, the alternating exponential series and
\eqref{eq:terminal-modal-full-optimum-x-upper} give
\[
 T=e^{-x}
 \ge1-x_0+\frac{x_0^2}{2}-\frac{x_0^3}{6}
       +\frac{x_0^4}{24}-\frac{x_0^5}{120}
 >\frac{2711}{3072}.
\]
Substitution of $2711/3072$ in
\eqref{eq:terminal-modal-proper-seed-optimum} leaves the exact slack
\[
 \frac18-
 \frac{2(2+2711/3072-3(2711/3072)^2)}
      {5(1+(2711/3072)^2)}
 =\frac{1469573}{671468200}>0.
\]
Thus $p_m^\star<q/8$ coordinatewise, and
\eqref{eq:terminal-modal-entry-f0-transport} proves the last bound.
\end{proof}

Let $p_k^{\rm lin}$ denote the full-face linear candidate obtained by omitting
the coordinate projection.  Its residual and global envelope are
\begin{equation}
 r_k^{\rm lin}=r(p_k^{\rm lin}),\qquad
 \delta_k^{\rm lin}=
 \max\left(0,\frac{\max_jr_{k,j}^{\rm lin}}{q^2}\right).
 \label{eq:terminal-modal-linear-residual-envelope}
\end{equation}
The linear error recurrence and \cref{lem:terminal-modal-propagators} give
unconditionally
\begin{equation}
 p_k^{\rm lin}-p^\star
 =\delta^k\left\{K_ke_0+J_k(qLf_0)\right\}.
 \label{eq:terminal-modal-position-candidate-kernel}
\end{equation}

\begin{lemma}[Late position dominates the global envelope; proved here]
\label{lem:terminal-modal-late-position-envelope}
Assume the static directed-remainder target
\begin{equation}
 \norm{u_+}_{1,D}\le\frac{21}{80}q^3.
 \label{eq:terminal-modal-static-target-assumed}
\end{equation}
If the literal execution agrees with the full-face linear candidate through
time $k-1$ and $qk\ge3/50$, then the time-$k$ raw candidate is strictly
positive and
\begin{equation}
 \min_jp_k^{\rm lin}(j)>\delta_k^{\rm lin}.
 \label{eq:terminal-modal-late-envelope-dominance}
\end{equation}
Consequently the literal execution agrees through time $k$, and its global
safe envelope is strictly unclipped at that time.
\end{lemma}

\begin{proof}
The kernels $K_k,L$, and $J_k$ are positive and have row masses $1,1$, and
$k$.  Hence \cref{lem:terminal-modal-entry-position-bounds} and
\eqref{eq:terminal-modal-position-candidate-kernel} give, with $s=qk$,
\begin{equation}
 \min_jp_k^{\rm lin}(j)
 >q\left\{\frac{1243}{160}
 -\delta^k\left(\frac{157}{20}+\frac{s}{8}\right)\right\}.
 \label{eq:terminal-modal-late-position-lower}
\end{equation}
On the other hand,
\cref{eq:terminal-modal-directed-positive-bound} and
\eqref{eq:terminal-modal-static-target-assumed} imply
\begin{equation}
 \delta_k^{\rm lin}
 \le\frac{21}{80}q\delta^k
       \left(1+\frac{k-1}{2m}\right)
 \le\frac{21}{80}q e^{-s}(1+8s).
 \label{eq:terminal-modal-late-envelope-upper}
\end{equation}
It therefore suffices that
\begin{equation}
 \frac{1243}{160}>
 e^{-s}\left(\frac{649}{80}+\frac{89}{40}s\right).
 \label{eq:terminal-modal-late-scalar-target}
\end{equation}
The right side is decreasing for $s\ge0$.  At $s=3/50$, use
$e^{-s}\le1-s+s^2/2$; the resulting exact slack is
\begin{equation}
 \frac{1243}{160}
 -\left(1-\frac{3}{50}+\frac12\left(\frac{3}{50}\right)^2\right)
  \left(\frac{649}{80}+\frac{89}{40}\frac{3}{50}\right)
 =\frac{1667}{625000}>0.
 \label{eq:terminal-modal-late-rational-slack}
\end{equation}
This proves \eqref{eq:terminal-modal-late-envelope-dominance}.
\end{proof}

The first full-face average residual also has the favorable sign.  This is
not an assumption on the open nonlinear regime; it is one last application
of the proper-prefix correction ledger.

\begin{lemma}[Full-entry average residual sign; proved here]
\label{lem:terminal-modal-first-full-average-sign}
The first full-face average residual satisfies
\begin{equation}
 \overline R_0
 =r\!\left(\frac{p_0+qv_0}{1+q}\right)<0
 \quad\hbox{coordinatewise}.
 \label{eq:terminal-modal-first-full-average-sign}
\end{equation}
\end{lemma}

\begin{proof}
First evaluate the zero-padded last proper raw point on a hypothetical
$(m+1)$-row prefix in which the new endpoint still has degree two.  Let
$\widetilde X_{m+1}$ be the negative residual of that point and define
$\widetilde K_{m+1}$ from it by the same ideal-correction split as in
\eqref{eq:terminal-modal-prefix-correction}.  On the old rows,
$\widetilde X_{m+1}=-r_0^x$.

The proof of \cref{lem:terminal-modal-finite-correction-sign} extends to this
one extra hypothetical transition.  Indeed, its last moving source is still
the standard degree-two source $n=m$; here $r=m-2$, $\chi^m\ge7/8$, and
$\delta^{m-1}\ge15/16$, so the scalar and variation estimates are unchanged.
For $N=m+1$, the initial and mass losses become exactly
\[
 \frac{3q}{5}+\frac{43q(m-1)}{75}
 =\frac{43}{1200}+\frac{1}{600m}.
\]
The derivative loss remains $1/90$, while the leading shared-coordinate
margin remains $19/320$.  Consequently
\begin{equation}
 \frac{\widetilde K_{m+1}(j)}{q^3\delta^{m-1}}
 \ge\frac{179}{14400}-\frac{1}{600m}
 \ge\frac{1429}{115200}>0,
 \qquad 0\le j<m.
 \label{eq:terminal-modal-hypothetical-final-correction}
\end{equation}

On the old rows the affine residual identity and optimum transport give
\begin{align}
 (1+q)\overline R_0(j)
 &=2r_0^x(j)-\delta A_m(j)\notag\\
 &=-2\left(\widetilde X_{m+1}(j)-\frac\delta2X_m(j)\right)\notag\\
 &=-2\left[
  G_{m+1}^{\rm id}(j)-\frac\delta2G_m^{\rm id}(j)
  +\widetilde K_{m+1}(j)
 \right]<0.
 \label{eq:terminal-modal-first-full-old-rows}
\end{align}
The ideal difference is zero at the seed and, by Pascal's identity, equals
$c_1G_m^{\rm id}(j)j/(m-j)\ge0$ for $1\le j<m$.

It remains to use the actual degree-one endpoint.  Since the new coordinate
was zero before the last proper step, exact center transport gives
\begin{align}
 y_0(m-1)&=
 \frac{2\sigma_m+q\{p^\star(m-1)-P_m\}}{1+q},
 &y_0(m)&=\frac{qp^\star(m)}{1+q},
 \label{eq:terminal-modal-first-full-endpoint-average}
\end{align}
where $y_0=(p_0+qv_0)/(1+q)$.  Substitution in the degree-one residual row,
followed by the full-optimum endpoint equation, yields
\begin{equation}
 (1+q)\overline R_0(m)
 =-2\eta\sigma_m+\eta qP_m+\frac{q^3}{5}.
 \label{eq:terminal-modal-first-full-endpoint-residual}
\end{equation}
Using \eqref{eq:terminal-modal-raw-frontier-lower},
$P_m/q^2\le33/16$, and $\eta\ge255/512$, the right side divided by $q^3$
is at most
\begin{equation}
 \frac{255}{512}\left(-2\frac{596861}{326570}+\frac{33}{16}\right)
 +\frac15
 =-\frac{46683733}{78684160}<0.
 \label{eq:terminal-modal-first-full-endpoint-margin}
\end{equation}
This proves the sign on every full-face row.
\end{proof}

We now isolate the remaining early target.  Put
\begin{equation}
 s_k^{\rm lin}=\delta^{-k}r_k^{\rm lin},
 \qquad X_k^{\rm lin}=-s_k^{\rm lin},
 \qquad
 \mathcal H_k=K_k-\frac12K_{k-1},
 \qquad
 \mathcal L_k=J_k-\frac12J_{k-1}.
 \label{eq:terminal-modal-half-retention-data}
\end{equation}
Here $K_{-1}$ is never used.  For $k\ge1$, $\mathcal H_k$ is positive with
row mass $1/2$: on the even cycle,
\[
 \mathcal H_k=\frac14\left(
  \mathcal S_+\mathcal A_+^{k-1}
 +\mathcal S_-\mathcal A_-^{k-1}\right).
\]
The kernel $\mathcal L_k$ is also positive, by the same factorization as
\cref{lem:terminal-modal-correction-positive-green} at $q=0$.

\begin{lemma}[Exact early half-retention interface; proved here]
\label{lem:terminal-modal-early-half-interface}
Suppose
\begin{equation}
 X_k^{\rm lin}-\frac12X_{k-1}^{\rm lin}\ge0
 \qquad(1\le k, qk<3/50).
 \label{eq:terminal-modal-early-half-retention}
\end{equation}
Then the literal projection is inactive, $r_k^x\le0$, and $\delta_k=0$
through this early window.

Moreover, with the directed ideal term
\begin{equation}
 \mathcal D_k=\mathcal A_+^kg^++\mathcal A_-^kg^-,
\end{equation}
the half-retention target is exactly
\begin{equation}
 X_k^{\rm lin}-\frac12X_{k-1}^{\rm lin}
 =-\left(\mathcal D_k-\frac12\mathcal D_{k-1}\right)
  -\mathcal H_kc-\mathcal L_ku.
 \label{eq:terminal-modal-early-half-kernel}
\end{equation}
The first term on the right is nonnegative, because
\[
 \mathcal D_k-\frac12\mathcal D_{k-1}
 =\frac12\mathcal S_+\mathcal A_+^{k-1}g^+
 +\frac12\mathcal S_-\mathcal A_-^{k-1}g^-\le0.
\]
Thus \eqref{eq:terminal-modal-early-half-kernel} is a finite
signed-convolution formulation of the remaining early nonlinear interface.
\end{lemma}

\begin{proof}
The proper chronology and positive final optimum displacement give
$v_0\ge p_0$, while the new endpoint has
$v_0(m)=p^\star(m)>0$.  Hence the first average point is strictly positive.
Together with \cref{lem:terminal-modal-first-full-average-sign}, this makes
the first raw update linear and gives
$r_1^{\rm lin}=B\overline R_0\le0$.  The literal first step therefore equals
the linear candidate.  For every later linear step, the affine
residual identity is
\begin{equation}
 r(y_k^{\rm lin})
 =\frac{2r_k^{\rm lin}-\delta r_{k-1}^{\rm lin}}{1+q}
 =-\frac{2\delta^k}{1+q}
  \left(X_k^{\rm lin}-\frac12X_{k-1}^{\rm lin}\right)\le0.
 \label{eq:terminal-modal-half-average-residual}
\end{equation}
Suppose inductively that the literal state equals the linear candidate through
time $k$.  If $r(y_k^{\rm lin})\le0$ and $v_k\ge p_k$, then the raw update
obeys $p_{k+1}^{\rm lin}\ge y_k^{\rm lin}\ge p_k^{\rm lin}>0$, so projection
is inactive and the literal next state again equals the candidate.  The center
formula gives $v_{k+1}\ge p_{k+1}$, while
$r_{k+1}^{\rm lin}=Br(y_k^{\rm lin})\le0$.  This proves the induction,
including $\delta_k=\delta_k^{\rm lin}=0$.

Finally substitute the directed decomposition from
\cref{lem:terminal-modal-directed-packet} into the definition of
$X_k^{\rm lin}$ and
subtract half the preceding time.  This gives
\eqref{eq:terminal-modal-early-half-kernel}; the displayed directed
difference follows from
$\mathcal A_\pm-\frac12I=\frac12\mathcal S_\pm$.
\end{proof}

Combining the lemmas gives a complete logical stitch: the static target
\eqref{eq:terminal-modal-static-target-assumed} and
\eqref{eq:terminal-modal-early-half-retention} through $qk<3/50$
imply the open regime lemma through the full horizon
\eqref{eq:terminal-modal-horizon}.  What remains unproved is the signed early
convolution, not any projection or envelope implication.  The sharpened entry
correction in \cref{lem:terminal-modal-entry-correction-margin} reduces this
shorter target to \eqref{eq:terminal-modal-short-early-convolution-open}.
That finite convolution bound remains open.  Extending the same separate
$q^3/16$ bound to the former endpoint $17/200$ is false on the deterministic
screen and is not used.  The coordinate-error monotonicity claim
that later face admissions can only improve $p-p_n^\star$ is false because
the restricted optimum itself changes at every admission.
